% Part: lambda-calculus
% Chapter: introduction
% Section: basic-pr-lambda

\documentclass[../../../include/open-logic-section]{subfiles}

\begin{document}

\olfileid{lam}{int}{bas}
\olsection{مُڈھلے بنیادی بازگشتی فنکشن \usetoken{S}{lambda definable} نیں}

\begin{lem}
فنکشن $\Zero$، $\Succ$ تے $\Proj{n}{i}$، !!{lambda definable} نیں۔
\end{lem}

\begin{proof}
$\Zero$ تاں بس
$\lambd[x][\lambd[y][y]]$ اے۔

جانشین فنکشن~$\Succ$ دی تعریف $\fn{Succ}(u) = \lambd[x][\lambd[y][x(uxy)]]$ اے۔
سوچو ایہ کیوں کم کردا اے۔ ہر چرچ عدد $\num{n}$ نوں وار وار عمل کرن
والا ٹرم سمجھو۔ کسے وی فنکشن~$f$ لئی $\fn{Succ}(\num{n},f)$
اوہ فنکشن اے جیہڑا انپٹ~$y$ لئی $f$ دا $n$ واری اطلاق کردا اے،
شروعات~$y$ توں کر کے۔ فیر اک واری ہور لاگو کردا اے۔

دِتّیاں دلیلاں وچوں اک دلیل چُنن والے فنکشوناں لئی کوئی ہور گل نہیں۔
$\fn{Proj}^n_i(x_0, \dots,
x_{n-1}) = x_i$۔ دوجے لفظاں وچ ساڈے لکھت دے رواج مطابق $\fn{Proj}^n_i$
لامبڈا ٹرم $\lambd[x_0][\dots \lambd[x_{n-1}][x_i]]$ اے۔
\end{proof}

\end{document}
